\begin{resume}
\ifisanon
    \ifisblind
        \isanonfalse
    \else
        \ifismaster
        该论文作者在学期间取得的阶段性成果（学术论文等）已满足我校硕士学位评阅相关要求。为避免阶段性成果信息对专家评阅学位论文本身造成干扰，特将论文作者的阶段性成果信息隐去。
        \else
        该论文作者在学期间取得的阶段性成果（学术论文等）已满足我校博士学位评阅相关要求。为避免阶段性成果信息对专家评阅学位论文本身造成干扰，特将论文作者的阶段性成果信息隐去。
        \fi
    \fi
\else\relax
\fi

\ifisanon\relax
\else
\ifisresumebib

\begingroup
    \settoggle{bbx:gbtype}{false}% 局部设置不输出文献类型和载体标识符
    \settoggle{bbx:gbannote}{true}% 局部设置输出注释信息
    \renewcommand{\bibfont}{\normalsize}
    \makeatletter
    \renewcommand*{\mkbibnamegiven}[1]{%
    \ifitemannotation{thesisauthor}
    {\ifbibliography{{\textbf{#1}}}{#1}}%
    {#1}\ifbibliography{\ifitemannotation{corresponding}{\textsuperscript{*}}{}}{}%
    }
    \renewcommand*{\mkbibnamefamily}[1]{%
    \ifitemannotation{thesisauthor}
    {\ifbibliography{{\textbf{#1}}}{#1}}
    {#1}}
    \def\blx@maxbibnames{9}
    \def\blx@minbibnames{2}
    \renewbibmacro*{author}{%
      \ifboolexpr{
        test \ifuseauthor
        and
        not test {\ifnameundef{author}}
      }
        {\ifisblind%
           \setcounter{numitemval}{1}%
           \whileboolexpr{test {\ifnumcomp{\value{numitemval}}{<}{\value{author}+1}}}
               {\ifitemannotation[author][default][\thenumitemval]{thesisauthor}%
                    {\setcounter{numitemslt}{\thenumitemval}}{}%
               \stepcounter{numitemval}%
               }%
           \ifnumcomp{\value{numitemslt}}{>}{0}
                {\textbf{\iffieldequalstr{userd}{chinese}{第\zhnum{numitemslt}作者}{\Ordinalstring{numitemslt}~Author}}, XXX.}
                {\printnames{author}}
         \else%
            \printnames{author}%
         \fi%%
         \iffieldundef{authortype}
           {}
           {\setunit{\printdelim{authortypedelim}}%
            \usebibmacro{authorstrg}}}
        {}}
    \makeatother

    \begin{refsection}[ref/resume.bib]
    \nocite{shu2024few,shu2025single}
    \printbibliography[heading=subbibliography,title={发表的学术论文}]
    \end{refsection}
\endgroup

\section*{专利成果}
\begin{enumerate}[label={[\arabic*]},labelsep=6pt,topsep=0pt,partopsep=0pt,
parsep=0pt,itemsep=1pt,leftmargin=0em,itemindent=42pt]
  \item 专利授权: “基于纯仿真数据训练的域泛化SAR目标识别方法、终端设备及存储介质”发明专利
  （专利号 ZL 2025 1 0009255.2，授权公告号 CN 119942267 B，申请日 2025年01月03日，授权公告日 2026年01月06日，排名2）
\end{enumerate}

\section*{竞赛获奖}
\begin{enumerate}[label={[\arabic*]},labelsep=6pt,topsep=0pt,partopsep=0pt,
parsep=0pt,itemsep=1pt,leftmargin=0em,itemindent=42pt]
  \item “华为杯”第二十一届中国研究生数学建模竞赛二等奖
  \item “兆易创新杯”第十九届中国研究生电子设计竞赛华中赛区二等奖
\end{enumerate}

\else

\section*{发表的学术论文}
\begin{enumerate}[label={[\arabic*]},labelsep=6pt,topsep=0pt,partopsep=0pt,
parsep=0pt,itemsep=1pt,leftmargin=0em,itemindent=42pt]
  \item Shu~W, Jiang~C, Zhan~R, et~al.
  Few-shot unsupervised domain adaptation for SAR target recognition based on input alignment and class confusion loss~[C].
  In 2024 IEEE International Conference on Signal, Information and Data Processing (ICSIDP).
  2024: 1--5.
  \item Shu~W, Zhan~R, Chen~S, et~al.
  Single-Domain Generalization via Multilevel Data Augmentation for SAR Target Recognition Training on Fully Simulated Data~[J].
  Remote Sensing.
  2025, 17~(17): 2966.
\end{enumerate}

\section*{专利成果}
\begin{enumerate}[label={[\arabic*]},labelsep=6pt,topsep=0pt,partopsep=0pt,
parsep=0pt,itemsep=1pt,leftmargin=0em,itemindent=42pt]
  \item 专利授权: “基于纯仿真数据训练的域泛化SAR目标识别方法、终端设备及存储介质”发明专利
  （专利号 ZL 2025 1 0009255.2，授权公告号 CN 119942267 B，申请日 2025年01月03日，授权公告日 2026年01月06日，排名2）
\end{enumerate}

\section*{竞赛获奖}
\begin{enumerate}[label={[\arabic*]},labelsep=6pt,topsep=0pt,partopsep=0pt,
parsep=0pt,itemsep=1pt,leftmargin=0em,itemindent=42pt]
  \item “华为杯”第二十一届中国研究生数学建模竞赛二等奖
  \item “兆易创新杯”第十九届中国研究生电子设计竞赛华中赛区二等奖
\end{enumerate}

\fi
\fi
\end{resume}
